\textbf{Phần 1. Trắc nghiệm lựa chọn}
\setcounterpageref{ex}{0}

\begin{ex} %[1D5V19]
Hình tam giác có bao nhiêu cạnh?
\choice
{$2$}
{$4$}
{\True $3$}
{$1$}
\loigiai{
Hình tam giác có đúng 3 cạnh.
\textbf{Đáp số:} $3$.\par
\textbf{Đáp án của câu hỏi là C}
}
\end{ex}

\begin{ex} %[1D5V19]
Hình vuông có bao nhiêu cạnh bằng nhau?
\choice
{\True $4$}
{$3$}
{$5$}
{$2$}
\loigiai{
Hình vuông có đúng 4 cạnh bằng nhau.
\textbf{Đáp số:} $4$.\par
\textbf{Đáp án của câu hỏi là A}
}
\end{ex}

\begin{ex} %[1D5V19]
Hình chữ nhật có bao nhiêu góc vuông?
\choice
{$3$}
{$5$}
{$2$}
{\True $4$}
\loigiai{
Hình chữ nhật có đúng 4 góc vuông.
\textbf{Đáp số:} $4$.\par
\textbf{Đáp án của câu hỏi là D}
}
\end{ex}

\begin{ex} %[1D5V19]
Khối lập phương có bao nhiêu mặt?
\choice
{$5$}
{\True $6$}
{$7$}
{$4$}
\loigiai{
Khối lập phương có đúng 6 mặt phẳng đều là hình vuông.
\textbf{Đáp số:} $6$.\par
\textbf{Đáp án của câu hỏi là B}
}
\end{ex}

\begin{ex} %[1D5V19]
Khối hộp chữ nhật có bao nhiêu đỉnh?
\choice
{$7$}
{$9$}
{\True $8$}
{$6$}
\loigiai{
Khối hộp chữ nhật có đúng 8 đỉnh.
\textbf{Đáp số:} $8$.\par
\textbf{Đáp án của câu hỏi là C}
}
\end{ex}

\begin{ex} %[1D5V19]
Tổng số cạnh của một hình tam giác và một hình chữ nhật là bao nhiêu?
\choice
{\True $7$}
{$6$}
{$8$}
{$5$}
\loigiai{
Hình tam giác có 3 cạnh, hình chữ nhật có 4 cạnh: $3 + 4 = 7$ (cạnh).
\textbf{Đáp số:} $7$.\par
\textbf{Đáp án của câu hỏi là A}
}
\end{ex}

\begin{ex} %[1D5V19]
Tổng số đỉnh của hai hình tam giác là bao nhiêu?
\choice
{$5$}
{$7$}
{$4$}
{\True $6$}
\loigiai{
Mỗi hình tam giác có 3 đỉnh. Hai hình tam giác có: $3 + 3 = 6$ (đỉnh).
\textbf{Đáp số:} $6$.\par
\textbf{Đáp án của câu hỏi là D}
}
\end{ex}

\begin{ex} %[1D5V19]
Một hình vuông có 4 cạnh. Nếu gấp đôi số cạnh đó lên thì được bao nhiêu cạnh?
\choice
{$7$}
{\True $8$}
{$9$}
{$6$}
\loigiai{
Gấp đôi số cạnh của hình vuông: $4 + 4 = 8$ (cạnh).
\textbf{Đáp số:} $8$.\par
\textbf{Đáp án của câu hỏi là B}
}
\end{ex}

\begin{ex} %[1D5V19]
Một hình vuông được chia đôi bởi một đường chéo tạo thành bao nhiêu hình tam giác?
\choice
{$1$}
{$3$}
{\True $2$}
{$4$}
\loigiai{
Đường chéo chia hình vuông thành $2$ hình tam giác bằng nhau.
\textbf{Đáp số:} $2$.\par
\textbf{Đáp án của câu hỏi là C}
}
\end{ex}

\begin{ex} %[1D5V19]
Cần ít nhất bao nhiêu que tính để xếp thành 2 hình tam giác dính liền chung 1 cạnh?
\choice
{\True $5$}
{$4$}
{$6$}
{$3$}
\loigiai{
Hình tam giác 1 cần 3 que, hình tam giác 2 chung 1 cạnh nên cần thêm 2 que: $3 + 2 = 5$ (que tính).
\textbf{Đáp số:} $5$.\par
\textbf{Đáp án của câu hỏi là A}
}
\end{ex}

\begin{ex} %[1D5V19]
Cần ít nhất bao nhiêu que tính để xếp thành 2 hình vuông dính liền chung 1 cạnh?
\choice
{$6$}
{$8$}
{$5$}
{\True $7$}
\loigiai{
Hình vuông 1 cần 4 que, hình vuông 2 chung 1 cạnh nên cần thêm 3 que: $4 + 3 = 7$ (que tính).
\textbf{Đáp số:} $7$.\par
\textbf{Đáp án của câu hỏi là D}
}
\end{ex}

\begin{ex} %[1D5V19]
Một mảnh giấy hình vuông có 4 góc. Cắt đi 1 góc theo đường chéo nhỏ. Hỏi mảnh giấy lớn còn lại có bao nhiêu góc?
\choice
{$4$}
{\True $5$}
{$6$}
{$3$}
\loigiai{
Số góc còn lại là: $4 - 1 + 2 = 5$ (góc).
\textbf{Đáp số:} $5$.\par
\textbf{Đáp án của câu hỏi là B}
}
\end{ex}

\textbf{Phần 2. Trắc nghiệm đúng sai}
\setcounterpageref{ex}{0}

\begin{ex} %[1D5V19]
Xét tính đúng hoặc sai của các khẳng định toán học sau đây:
\choiceTF
{\True Xếp 4 khối lập phương nhỏ thành một khối lập phương lớn hơn có 2 tầng. mỗi tầng có $4$ khối lập phương nhỏ.}
{Trong một bức tranh có: 2 hình tròn, 3 hình vuông và 4 hình tam giác. có tất cả $8$ hình phẳng.}
{\True Có 10 que tính bằng nhau. Bạn Nam xếp được 2 hình vuông riêng biệt. Nam còn thừa $2$ que tính.}
{Có 10 que tính bằng nhau. Bạn Mai xếp được 3 hình tam giác riêng biệt. Mai còn thừa $0$ que tính.}
\loigiai{
\begin{itemize}
    \item Ý a) Đúng: Một khối lập phương lớn $2 \times 2 \times 2$ có 8 khối, mỗi tầng có $4$ khối.
\textbf{Đáp số:} $4$.
    \item Ý b) Sai: Kết quả đúng là $9$, không phải $8$. Tổng số hình phẳng là: $2 + 3 + 4 = 9$ (hình).
\textbf{Đáp số:} $9$.
    \item Ý c) Đúng: Xếp 2 hình vuông dùng: $4 + 4 = 8$ (que tính).\nSố que tính còn thừa là: $10 - 8 = 2$ (que tính).
\textbf{Đáp số:} $2$.
    \item Ý d) Sai: Kết quả đúng là $1$, không phải $0$. Xếp 3 hình tam giác dùng: $3 + 3 + 3 = 9$ (que tính).\nSố que tính còn thừa là: $10 - 9 = 1$ (que tính).
\textbf{Đáp số:} $1$.
\end{itemize}
\begin{itemize}
    \item a) \textbf{Đúng}.
    \item b) \textbf{Sai}.
    \item c) \textbf{Đúng}.
    \item d) \textbf{Sai}.
\end{itemize}
}
\end{ex}

\begin{ex} %[1D5V19]
Xét tính đúng hoặc sai của các bài toán và phép tính sau đây:
\choiceTF
{Kết quả của bài toán «Khối rubik tiêu chuẩn có dạng hình khối gì? (Khối lập phương: ghi 1; Khối hộp chữ nhật: ghi 2)» là $0$.}
{\True Kết quả của bài toán «Chiếc hộp đựng bánh pizza lớn thường có dạng khối gì? (Khối hộp chữ nhật dẹt: ghi 2; Khối lập phương: ghi 1)» là $2$.}
{Kết quả của bài toán «Một hình chữ nhật được ghép từ 2 hình vuông bằng nhau có cạnh dài 3 cm. Chiều dài của hình chữ nhật đó là bao nhiêu xăng-ti-mét» là $7\text{ cm}$.}
{\True Có 10 que tính bằng nhau. Bạn Nam xếp được 2 hình vuông riêng biệt. Nam còn thừa $2$ que tính.}
\loigiai{
\begin{itemize}
    \item Ý a) Sai: Kết quả đúng là $1$, không phải $0$. Khối rubik có dạng khối lập phương (ghi số 1).
\textbf{Đáp số:} $1$.
    \item Ý b) Đúng: Hộp bánh pizza dẹt có dạng khối hộp chữ nhật (ghi số 2).
\textbf{Đáp số:} $2$.
    \item Ý c) Sai: Kết quả đúng là $6$, không phải $7$. Chiều dài hình chữ nhật là: $3 + 3 = 6$ (cm).
\textbf{Đáp số:} $6$.
    \item Ý d) Đúng: Xếp 2 hình vuông dùng: $4 + 4 = 8$ (que tính).\nSố que tính còn thừa là: $10 - 8 = 2$ (que tính).
\textbf{Đáp số:} $2$.
\end{itemize}
\begin{itemize}
    \item a) \textbf{Sai}.
    \item b) \textbf{Đúng}.
    \item c) \textbf{Sai}.
    \item d) \textbf{Đúng}.
\end{itemize}
}
\end{ex}

\textbf{Phần 3. Trả lời ngắn}
\setcounterpageref{ex}{0}

\begin{ex} %[1D5V19]
Tổng số mặt của một khối lập phương và một khối hộp chữ nhật là bao nhiêu mặt?
\shortans{12}
\loigiai{
\textbf{Đáp số:} $12$\par
Mỗi khối có 6 mặt: $6 + 6 = 12$ (mặt).
\textbf{Đáp số:} $12$.
}
\end{ex}

\begin{ex} %[1D5V19]
Bé dùng 6 que tính xếp được 1 hình chữ nhật có chiều dài 2 que, chiều rộng 1 que. Hỏi chu vi hình chữ nhật đó gồm bao nhiêu que tính?
\shortans{6}
\loigiai{
\textbf{Đáp số:} $6$\par
Chu vi hình chữ nhật đó gồm: $2 + 2 + 1 + 1 = 6$ (que tính).
\textbf{Đáp số:} $6$.
}
\end{ex}

\begin{ex} %[1D5V19]
Một hình tam giác có độ dài các cạnh là 3 cm, 4 cm và 3 cm. Tổng độ dài ba cạnh của hình tam giác đó là bao nhiêu xăng-ti-mét?
\shortans{10}
\loigiai{
\textbf{Đáp số:} $10$\par
Tổng độ dài ba cạnh là: $3 + 4 + 3 = 10$ (cm).
\textbf{Đáp số:} $10$.
}
\end{ex}

\begin{ex} %[1D5V19]
Khối hình học nào không thể lăn tròn được trên mặt bàn? (Khối cầu: 1; Khối lập phương: 2).
\shortans{2}
\loigiai{
\textbf{Đáp số:} $2$\par
Khối lập phương có các mặt phẳng nên không lăn trơn được như khối cầu (ghi số 2).
\textbf{Đáp số:} $2$.
}
\end{ex}

\begin{ex} %[1D5V19]
Trong các hình: hình tam giác, hình vuông, hình tròn, hình chữ nhật; hình nào không có đoạn thẳng nào bao quanh? (Hình tam giác: 1; Hình tròn: 2; Hình vuông: 3).
\shortans{2}
\loigiai{
\textbf{Đáp số:} $2$\par
Hình tròn được bao quanh bởi đường cong khép kín, không có đoạn thẳng (ghi số 2).
\textbf{Đáp số:} $2$.
}
\end{ex}

\begin{ex} %[1D5V19]
Một hình vuông có chu vi gồm 8 que tính. Hỏi mỗi cạnh của hình vuông đó gồm bao nhiêu que tính?
\shortans{2}
\loigiai{
\textbf{Đáp số:} $2$\par
Vì hình vuông có 4 cạnh bằng nhau và $2 + 2 + 2 + 2 = 8$ nên mỗi cạnh gồm $2$ que tính.
\textbf{Đáp số:} $2$.
}
\end{ex}
